\documentclass[12pt]{article}

\include{../style}

\begin{document}

\title{Cosc 30: Discrete Mathematics}

\author{Prishita Dharampal}
\date{}
\maketitle
\vspace{-0.3in}


\begin{problem}{1}
Prove that a graph $G$ is a tree if and only if $G$ is (C) connected and (MinC) deleting an edge from $G$ disconnects the graph.
\end{problem}

\begin{solution}
    \bbni

    Graph $G$ is a tree if and only if $G$ is connected and acyclic. So we just need to prove that the edge deleting property (MinC) is the same as $G$ being acyclic. 

    \bbni 
    $(\implies)$ Assume that $G$ is connected and acyclic. We want to show that then it is also minimally connected. 

    Consider two arbitrary vertices $u$ and $v$ in $G$. Since $G$ is connected there exist at least one path between them, and since $G$ is acyclic there exist at most one path between them. That is, there is exactly one path between $u$ and $v$, and deleting an arbitrary edge on this path would seperate the graph into two disjoint connected components. 

    \bbni
    $(\impliedby)$ Assume $G$ is connected and minimally connected. We want to show that then it is acyclic. 

    For the sake of contradiction, assume that $G$ contains a cyclic. Now consider two adjacent vertices $u$ and $v$ on this cycle. Since the cycle provides another path to get from $u$ to $v$, removing the edge between the two vertices would still leave $G$ connected, contradicting the assumption that $G$ is minimally connected. Hence, $G$ cannot contain a cycle. 

    \bbni
    Hence, a graph is a tree if and only if it is connected and deleting an edge from the graph disconnects it. 
\end{solution}

\newpage 

\begin{problem}{2}
    Prove that if an $n$-node graph $G$ has (E) exactly $n - 1$ edges and is (C) connected, then $G$ is (A) acyclic (and thus a tree).
\end{problem}

\begin{solution}
    \bbni
    
    Assume $G$ has $n$ nodes, exactly $n-1$ edges and is connected. We want to show that $G$ is acyclic (and thus a tree). 
    
    Assume for the sake of contradiction that $G$ contains a cycle. Let this cycle have $k$ vertices. Since a cycle has the same number of edges as vertices, the cycle contains $k$ edges, where $k \ge 2$.    
    
    These $k$ vertices already form a connected component using $k$ edges. We are left with $n-k$ vertices outside the cycle. In order for $G$ to be connected, each of these $n-k$ vertices must be connected to the cycle (or to vertices already attached to it). Attaching each new vertex to an existing connected component requires at least one edge.
    
    Therefore, we need at least $n-k$ additional edges to connect the remaining vertices to the cycle. Altogether, the total number of edges in $G$ must be at least
    \[k + (n-k) = n\]
    So, we have $n-1-k$ edges, hence this gives us a contradiction. Therefore $G$ must be acyclic, and thus is a tree.
\end{solution}
\end{document}